% ============================================================================
% graficas/rango.tex
% Gráfica de Rango SUS: barra de 0 a 100, línea de referencia 68 y rango mín-máx
% Colores se cargan en el preámbulo principal (graficas/colores.tex)
% ============================================================================

\begin{tikzpicture}[x=0.095cm, y=1cm]
  % Usar macros definidas por sus.definir_macros() (valores con punto decimal)
  \pgfmathsetmacro{\minRaw}{\susMinimoRaw}
  \pgfmathsetmacro{\maxRaw}{\susMaximoRaw}
  \pgfmathsetmacro{\mediaRaw}{\susMediaRaw}

  % Fondo de aceptabilidad en la barra horizontal (y de 0 a 0.7)
  \fill[susRojo!25] (0,0) rectangle (50,0.7);
  \fill[susAmarillo!25] (50,0) rectangle (70,0.7);
  \fill[susVerde!25] (70,0) rectangle (100,0.7);

  % Borde exterior de la barra
  \draw[thick, gray!70] (0,0) rectangle (100,0.7);

  % Marcas de división en el eje
  \foreach \x in {0, 20, 40, 60, 80, 100} {
    \draw[gray!70] (\x, 0) -- (\x, -0.15) node[below, font=\scriptsize] {\x};
  }

  % Línea de referencia estándar de la industria (SUS = 68)
  \draw[susReferencia, very thick, dashed] (68, -0.3) -- (68, 1.3);
  \node[susReferencia, above, font=\scriptsize\bfseries, align=center] at (68, 1.35) {Media global SUS\\(68)};

  % Intervalo Mínimo - Máximo de los participantes
  \fill[susPrincipal!40, rounded corners=2pt] (\minRaw, 0.15) rectangle (\maxRaw, 0.55);
  \draw[susPrincipal, thick] (\minRaw, 0.35) -- (\maxRaw, 0.35);

  % Extremo mínimo
  \draw[susPrincipal, line width=1.2pt] (\minRaw, 0.15) -- (\minRaw, 0.55);
  \node[susPrincipal, below, font=\scriptsize\bfseries] at (\minRaw, -0.45) {Mín: \susMinimo};
  \draw[->, >=stealth, susPrincipal, thin] (\minRaw, -0.42) -- (\minRaw, 0.1);

  % Extremo máximo
  \draw[susPrincipal, line width=1.2pt] (\maxRaw, 0.15) -- (\maxRaw, 0.55);
  \node[susPrincipal, below, font=\scriptsize\bfseries] at (\maxRaw, -0.45) {Máx: \susMaximo};
  \draw[->, >=stealth, susPrincipal, thin] (\maxRaw, -0.42) -- (\maxRaw, 0.1);

  % Marcador de la media muestral obtenida
  \fill[susPrincipal] (\mediaRaw, 0.35) circle (3pt);
  \node[susPrincipal, above, font=\small\bfseries] at (\mediaRaw, 0.75) {Media: \susMedia};
  \draw[->, >=stealth, susPrincipal, thick] (\mediaRaw, 0.72) -- (\mediaRaw, 0.45);
\end{tikzpicture}
